\subsection{Constructed continuations and attained qualification}\label{sec:study73}
The readout theorem is executed in four predeclared backward constructions in the original two-state, two-date, scalar-control economy. Two midpoint covers, with eight or sixteen cells per coordinate, are crossed with additive shallow ReLU and degree-four Bernstein features. Each basis has nine coefficients bounded by eight. At each date the future is the continuation actually constructed at the next date. Exact conditional moments supply thirty-three continuous feasible action values per node and thirty-three robust, downward-quantized acquired-action proposals. The original state-dependent capacity and continuous innovation law are unchanged.

\begin{table}[htbp]
\centering
\caption{Attained continuation qualification in the original investment economy}\label{tab:qualification73}
{\small\setlength{\tabcolsep}{4pt}
\begin{tabular}{rlrrr}
\hline
$n$ & Basis & $\Omega_h$ upper & Regret upper & Max. dual gap \\
\hline
8 & ReLU & 5.39764 & 8.46716 & 0.000000014 \\
8 & Bernstein & 5.36565 & 8.11941 & 0.000000000 \\
16 & ReLU & 3.73299 & 5.28939 & 0.000000014 \\
16 & Bernstein & 4.43474 & 5.81890 & 0.000000000 \\
\hline
\end{tabular}}
\par\smallskip
{\footnotesize Each row is one predeclared, complete two-date construction. Bounds are rounded upward; the exact fractions and dates are retained. Regret upper denotes the proved original-optimum account, not an estimated loss or an attained economic equivalence margin. The maximum dual gap concerns only the declared fixed-feature, fixed-future readout objectives.}
\end{table}


The table reports the attained two-date residual account $\Omega_h$, the original-optimum bound $\Omega_h+H_h^\pi$, and the largest exact one-step optimization gap. A smaller qualification objective is not itself an identified expected-cost gain. The matched least-squares ablation in the appendix uses each service's identical features and fixed-future labels; it tests the fitting objective, not a separately propagated controller. Both fitting and uniform qualification are paid, and the explicit state cover is tensor based.

All eight datewise optimization gaps are below $1.4\times10^{-8}$, against the predeclared tolerance $1/65536$. Their qualified residual bounds are strictly below the matched label-loss bounds; the reductions range from about 18 to 63 percent. The resulting global regret bounds remain conservative, as the table makes explicit. Those objective reductions do not establish a narrow economic equivalence interval or cheaper qualification at the earlier $1/256$ margin.

Every new action value is independently reconstructed using the retained original rational terminal formula or a different exact integration method for the fitted continuation. Continuous-action lower covers, whole-cell choices and primal--dual bounds are then recomputed. These constructions add actual qualified continuations, not new independent cost observations. The stopping theorem has separate exact finite-model regressions; the earlier fixed-sample equivalence intervals are not retroactively given an anytime interpretation.
